\documentclass[10pt,a4paper]{amsart}
\usepackage{amsmath,amssymb,amsthm,graphicx}
\usepackage{verbatim}
\newenvironment{reference}{\comment}{\endcomment}
%\newenvironment{reference}{}{}
\newenvironment{slogan}{\comment}{\endcomment}
\newenvironment{history}{\comment}{\endcomment}
% For commutative diagrams we use Xy-pic
\usepackage[all]{xy}

% We use 2cell for 2-commutative diagrams.
\xyoption{2cell}
% Theorem environments.
%
\theoremstyle{plain}
\newtheorem{theorem}[subsection]{Theorem}
\newtheorem{proposition}[subsection]{Proposition}
\newtheorem{lemma}[subsection]{Lemma}

\theoremstyle{definition}
\newtheorem{definition}[subsection]{Definition}
\newtheorem{example}[subsection]{Example}
\newtheorem{exercise}[subsection]{Exercise}
\newtheorem{situation}[subsection]{Situation}

\theoremstyle{remark}
\newtheorem{remark}[subsection]{Remark}
\newtheorem{remarks}[subsection]{Remarks}

\numberwithin{equation}{subsection}

% Macros
%
\def\lim{\mathop{\mathrm{lim}}\nolimits}
\def\colim{\mathop{\mathrm{colim}}\nolimits}
\def\Spec{\mathop{\mathrm{Spec}}}
\def\Hom{\mathop{\mathrm{Hom}}\nolimits}
\def\Ext{\mathop{\mathrm{Ext}}\nolimits}
\def\SheafHom{\mathop{\mathcal{H}\!\mathit{om}}\nolimits}
\def\SheafExt{\mathop{\mathcal{E}\!\mathit{xt}}\nolimits}
\def\Sch{\mathit{Sch}}
\def\Mor{\mathop{\mathrm{Mor}}\nolimits}
\def\Ob{\mathop{\mathrm{Ob}}\nolimits}
\def\Sh{\mathop{\mathit{Sh}}\nolimits}
\def\NL{\mathop{N\!L}\nolimits}
\def\CH{\mathop{\mathrm{CH}}\nolimits}
\def\proetale{{pro\text{-}\acute{e}tale}}
\def\etale{{\acute{e}tale}}
\def\QCoh{\mathit{QCoh}}
\def\Ker{\mathop{\mathrm{Ker}}}
\def\Im{\mathop{\mathrm{Im}}}
\def\Coker{\mathop{\mathrm{Coker}}}
\def\Coim{\mathop{\mathrm{Coim}}}

% Boxtimes
%
\DeclareMathSymbol{\boxtimes}{\mathbin}{AMSa}{"02}

%
% Macros for moduli stacks/spaces
%
\def\QCohstack{\mathcal{QC}\!\mathit{oh}}
\def\Cohstack{\mathcal{C}\!\mathit{oh}}
\def\Spacesstack{\mathcal{S}\!\mathit{paces}}
\def\Quotfunctor{\mathrm{Quot}}
\def\Hilbfunctor{\mathrm{Hilb}}
\def\Curvesstack{\mathcal{C}\!\mathit{urves}}
\def\Polarizedstack{\mathcal{P}\!\mathit{olarized}}
\def\Complexesstack{\mathcal{C}\!\mathit{omplexes}}
% \Pic is the operator that assigns to X its picard group, usage \Pic(X)
% \Picardstack_{X/B} denotes the Picard stack of X over B
% \Picardfunctor_{X/B} denotes the Picard functor of X over B
\def\Pic{\mathop{\mathrm{Pic}}\nolimits}
\def\Picardstack{\mathcal{P}\!\mathit{ic}}
\def\Picardfunctor{\mathrm{Pic}}
\def\Deformationcategory{\mathcal{D}\!\mathit{ef}}
\usepackage{cleveref}
\crefname{originalreference}{original reference}{original references}
\Crefname{originalreference}{Original reference}{Original references}
\makeatletter
\newlabel{algebra-lemma-ass-zero}{{algebra-lemma-ass-zero}{source}{}{sourcepointer-1}{}}
\newlabel{algebra-lemma-ass-zero@cref}{{[originalreference][0][]algebra-lemma-ass-zero}{[1][0][]source}{}{}{}}
\makeatother
\crefname{originaldefinition}{original definition}{original definitions}
\Crefname{originaldefinition}{Original definition}{Original definitions}
\makeatletter
\newlabel{definition-depth-complex}{{definition-depth-complex}{source}{}{sourcepointer-2}{}}
\newlabel{definition-depth-complex@cref}{{[originaldefinition][0][]definition-depth-complex}{[1][0][]source}{}{}{}}
\makeatother
\crefname{originalreference}{original reference}{original references}
\Crefname{originalreference}{Original reference}{Original references}
\makeatletter
\newlabel{dualizing-lemma-special-case-local-duality}{{dualizing-lemma-special-case-local-duality}{source}{}{sourcepointer-3}{}}
\newlabel{dualizing-lemma-special-case-local-duality@cref}{{[originalreference][0][]dualizing-lemma-special-case-local-duality}{[1][0][]source}{}{}{}}
\makeatother
\crefname{originalreference}{original reference}{original references}
\Crefname{originalreference}{Original reference}{Original references}
\makeatletter
\newlabel{dualizing-section-dimension-function}{{dualizing-section-dimension-function}{source}{}{sourcepointer-4}{}}
\newlabel{dualizing-section-dimension-function@cref}{{[originalreference][0][]dualizing-section-dimension-function}{[1][0][]source}{}{}{}}
\makeatother
\crefname{originallemma}{original lemma}{original lemmas}
\Crefname{originallemma}{Original lemma}{Original lemmas}
\makeatletter
\newlabel{lemma-local-cohomology-ss}{{lemma-local-cohomology-ss}{source}{}{sourcepointer-5}{}}
\newlabel{lemma-local-cohomology-ss@cref}{{[originallemma][0][]lemma-local-cohomology-ss}{[1][0][]source}{}{}{}}
\makeatother
\crefname{originallemma}{original lemma}{original lemmas}
\Crefname{originallemma}{Original lemma}{Original lemmas}
\makeatletter
\newlabel{lemma-torsion-change-rings}{{lemma-torsion-change-rings}{source}{}{sourcepointer-6}{}}
\newlabel{lemma-torsion-change-rings@cref}{{[originallemma][0][]lemma-torsion-change-rings}{[1][0][]source}{}{}{}}
\makeatother
\crefname{originalproposition}{original proposition}{original propositions}
\Crefname{originalproposition}{Original proposition}{Original propositions}
\makeatletter
\newlabel{proposition-annihilator}{{proposition-annihilator}{source}{}{sourcepointer-7}{}}
\newlabel{proposition-annihilator@cref}{{[originalproposition][0][]proposition-annihilator}{[1][0][]source}{}{}{}}
\makeatother
\crefname{originalreference}{original reference}{original references}
\Crefname{originalreference}{Original reference}{Original references}
\makeatletter
\newlabel{topology-lemma-Noetherian}{{topology-lemma-Noetherian}{source}{}{sourcepointer-8}{}}
\newlabel{topology-lemma-Noetherian@cref}{{[originalreference][0][]topology-lemma-Noetherian}{[1][0][]source}{}{}{}}
\makeatother
\title{Annihilation of bounded local cohomology of a coherent complex by an ideal supported away from a specialization-stable subset}
\author{The Stacks Project Authors}
\date{stacks-project:local-cohomology.tex:local-snapshot-20261012; GNU FDL 1.2 or later}
\begin{document}
\maketitle
\noindent\textit{Editorial note.} All delivered mathematical text is the author's own, copied from the pinned Stacks Project chapter source sources/stacks-local-cohomology/pinned-original-source.tex; no statement or proof step was rewritten, reordered or removed.\par\medskip
\noindent\textit{Editorial note.} Preamble provenance: the Stacks Project keeps one single shared preamble for every chapter; this card takes it directly from the staged sources/stacks-local-cohomology/preamble.tex as three verbatim spans (lines 8-12, 14-18 and 154-217), so no derived file was written and nothing was added or redefined. The wrapper class and the base packages are supplied by the assembler.\par\medskip
\noindent\textit{Editorial note.} Bibliography provenance: neither the counted statement nor its proof contains a citation (this lane rescans the whole span before writing the card), so the card carries no bibliography and no entry was reconstructed.\par\medskip
\noindent\textit{Editorial note.} Selection provenance: this card was opened under the file-based control contract (CONTROL.md), with the mechanical pre-card gate run on its statement and proof interval before the card was written.\par\medskip
\noindent\textit{Allowed standard prerequisites (editorial).}\begin{itemize}
\item Dualizing complexes, their dimension functions and the local duality spectral sequence of another chapter. Set up in the cited chapter's section on dimension functions and invoked inside the counted proof at several places (cross-chapter pointers). Supplies the dualizing object, the dimension function and the spectral sequence on which the whole argument runs.
\item The earlier annihilator proposition that the counted statement generalises, and the notion of the depth of a complex at a prime. The author's own proposition labelled proposition-annihilator and the cross-chapter definition of depth (pointers). Provide the model proof and the invariant appearing in the second condition.
\end{itemize}\medskip
\ifdefined\hypertarget\hypertarget{sourcepointer-1}{}\fi
\noindent\textit{Source reference.} Original reference algebra-lemma-ass-zero: sources/stacks-local-cohomology/pinned-original-source.tex, source lines 2323--2323. Navigation pointer only: the counted text refers to the Stacks item algebra-lemma-ass-zero, which belongs to another chapter of the project and is not part of this card. The pointer carries the target's own label as its number because this chapter's staged snapshot does not contain the other chapter's tag table; its content is not reproduced here.\par\medskip
\ifdefined\hypertarget\hypertarget{sourcepointer-2}{}\fi
\noindent\textit{Source reference.} Original definition definition-depth-complex: sources/stacks-local-cohomology/pinned-original-source.tex, source lines 3138--3138. Navigation pointer only: the counted text refers to the Stacks item definition-depth-complex, which belongs to another chapter of the project and is not part of this card. The pointer carries the target's own label as its number because this chapter's staged snapshot does not contain the other chapter's tag table; its content is not reproduced here.\par\medskip
\ifdefined\hypertarget\hypertarget{sourcepointer-3}{}\fi
\noindent\textit{Source reference.} Original reference dualizing-lemma-special-case-local-duality: sources/stacks-local-cohomology/pinned-original-source.tex, source lines 844--844. Navigation pointer only: the counted text refers to the Stacks item dualizing-lemma-special-case-local-duality, which belongs to another chapter of the project and is not part of this card. The pointer carries the target's own label as its number because this chapter's staged snapshot does not contain the other chapter's tag table; its content is not reproduced here.\par\medskip
\ifdefined\hypertarget\hypertarget{sourcepointer-4}{}\fi
\noindent\textit{Source reference.} Original reference dualizing-section-dimension-function: sources/stacks-local-cohomology/pinned-original-source.tex, source lines 2276--2276. Navigation pointer only: the counted text refers to the Stacks item dualizing-section-dimension-function, which belongs to another chapter of the project and is not part of this card. The pointer carries the target's own label as its number because this chapter's staged snapshot does not contain the other chapter's tag table; its content is not reproduced here.\par\medskip
\ifdefined\hypertarget\hypertarget{sourcepointer-5}{}\fi
\noindent\textit{Source reference.} Original lemma lemma-local-cohomology-ss: sources/stacks-local-cohomology/pinned-original-source.tex, source lines 1214--1214. Navigation pointer only: the counted text refers to the Stacks item lemma-local-cohomology-ss, which belongs to another chapter of the project and is not part of this card. The pointer carries the target's own label as its number because this chapter's staged snapshot does not contain the other chapter's tag table; its content is not reproduced here.\par\medskip
\ifdefined\hypertarget\hypertarget{sourcepointer-6}{}\fi
\noindent\textit{Source reference.} Original lemma lemma-torsion-change-rings: sources/stacks-local-cohomology/pinned-original-source.tex, source lines 2333--2333. Navigation pointer only: the counted text refers to the Stacks item lemma-torsion-change-rings, which belongs to another chapter of the project and is not part of this card. The pointer carries the target's own label as its number because this chapter's staged snapshot does not contain the other chapter's tag table; its content is not reproduced here.\par\medskip
\ifdefined\hypertarget\hypertarget{sourcepointer-7}{}\fi
\noindent\textit{Source reference.} Original proposition proposition-annihilator: sources/stacks-local-cohomology/pinned-original-source.tex, source lines 2454--2454. Navigation pointer only: the counted text refers to the Stacks item proposition-annihilator, which belongs to another chapter of the project and is not part of this card. The pointer carries the target's own label as its number because this chapter's staged snapshot does not contain the other chapter's tag table; its content is not reproduced here.\par\medskip
\ifdefined\hypertarget\hypertarget{sourcepointer-8}{}\fi
\noindent\textit{Source reference.} Original reference topology-lemma-Noetherian: sources/stacks-local-cohomology/pinned-original-source.tex, source lines 2377--2377. Navigation pointer only: the counted text refers to the Stacks item topology-lemma-Noetherian, which belongs to another chapter of the project and is not part of this card. The pointer carries the target's own label as its number because this chapter's staged snapshot does not contain the other chapter's tag table; its content is not reproduced here.\par\medskip
\begin{proposition}
\label{proposition-annihilator-complex}
Let $A$ be a Noetherian ring which has a dualizing complex.
Let $T \subset T' \subset \Spec(A)$ be subsets stable under
specialization. Let $s \in \mathbf{Z}$. Let $K$ be an object of
$D_{\textit{Coh}}^+(A)$. The following are equivalent
\begin{enumerate}
\item there exists an ideal $J \subset A$ with $V(J) \subset T'$
such that $J$ annihilates $H^i_T(K)$ for $i \leq s$, and
\item for all $\mathfrak p \not \in T'$,
$\mathfrak q \in T$ with $\mathfrak p \subset \mathfrak q$
we have
$$
\text{depth}_{A_\mathfrak p}(K_\mathfrak p) +
\dim((A/\mathfrak p)_\mathfrak q) > s
$$
\end{enumerate}
\end{proposition}

\begin{proof}
This lemma is the natural generalization of
Proposition \ref{proposition-annihilator}
whose proof the reader should read first.
Let $\omega_A^\bullet$ be a dualizing complex. Let $\delta$ be its
dimension function, see Dualizing Complexes, Section
\ref{dualizing-section-dimension-function}.
An important role will be played by the finite $A$-modules
$$
E^i = \Ext_A^i(K, \omega_A^\bullet)
$$
For $\mathfrak p \subset A$ we will write $H^i_\mathfrak p$ to denote the
local cohomology of an object of $D(A_\mathfrak p)$ with respect to
$\mathfrak pA_\mathfrak p$. Then we see that
the $\mathfrak pA_\mathfrak p$-adic completion of
$$
(E^i)_\mathfrak p =
\Ext^{\delta(\mathfrak p) + i}_{A_\mathfrak p}(K_\mathfrak p,
(\omega_A^\bullet)_\mathfrak p[-\delta(\mathfrak p)])
$$
is Matlis dual to
$$
H^{-\delta(\mathfrak p) - i}_{\mathfrak p}(K_\mathfrak p)
$$
by
Dualizing Complexes, Lemma \ref{dualizing-lemma-special-case-local-duality}.
In particular we deduce from this the
following fact: an ideal $J \subset A$ annihilates
$(E^i)_\mathfrak p$ if and only if $J$ annihilates
$H^{-\delta(\mathfrak p) - i}_{\mathfrak p}(K_\mathfrak p)$.

\medskip\noindent
Set $T_n = \{\mathfrak p \in T \mid \delta(\mathfrak p) \leq n\}$.
As $\delta$ is a bounded function, we see that
$T_a = \emptyset$ for $a \ll 0$ and $T_b = T$ for $b \gg 0$.

\medskip\noindent
Assume (2). Let us prove the existence of $J$ as in (1).
We will use a double induction to do this. For $i \leq s$
consider the induction hypothesis $IH_i$:
$H^a_T(K)$ is annihilated by some $J \subset A$ with $V(J) \subset T'$
for $a \leq i$. The case $IH_i$ is trivial for $i$ small
enough because $K$ is bounded below.

\medskip\noindent
Induction step. Assume $IH_{i - 1}$ holds for some $i \leq s$.
Pick $J'$ with $V(J') \subset T'$ annihilating $H^a_T(K)$ for
$a \leq i - 1$ (the induction hypothesis guarantees we can
do this). We will show by descending induction on $n$
that there exists an ideal $J$ with $V(J) \subset T'$ such that the
associated primes of $J H^i_T(K)$ are in $T_n$.
For $n \ll 0$ this implies $JH^i_T(K) = 0$ 
(Algebra, Lemma \ref{algebra-lemma-ass-zero})
and hence $IH_i$ will hold.
The base case $n \gg 0$ is trivial because $T = T_n$ in this case
and all associated primes of $H^i_T(K)$ are in $T$.

\medskip\noindent
Thus we assume given $J$ with the property for $n$.
Let $\mathfrak q \in T_n$. Let $T_\mathfrak q \subset \Spec(A_\mathfrak q)$
be the inverse image of $T$. We have
$H^j_T(K)_\mathfrak q = H^j_{T_\mathfrak q}(K_\mathfrak q)$
by Lemma \ref{lemma-torsion-change-rings}.
Consider the spectral sequence
$$
H_\mathfrak q^p(H^q_{T_\mathfrak q}(K_\mathfrak q))
\Rightarrow
H^{p + q}_\mathfrak q(K_\mathfrak q)
$$
of Lemma \ref{lemma-local-cohomology-ss}.
Below we will find an ideal $J'' \subset A$ with $V(J'') \subset T'$
such that $H^i_\mathfrak q(K_\mathfrak q)$ is annihilated by $J''$ for all
$\mathfrak q \in T_n \setminus T_{n - 1}$.
Claim: $J (J')^i J''$ will work for $n - 1$.
Namely, let $\mathfrak q \in T_n \setminus T_{n - 1}$.
The spectral sequence above defines a filtration
$$
E_\infty^{0, i} = E_{i + 2}^{0, i} \subset \ldots \subset E_3^{0, i} \subset
E_2^{0, i} = H^0_\mathfrak q(H^i_{T_\mathfrak q}(K_\mathfrak q))
$$
The module $E_\infty^{0, i}$ is annihilated by $J''$.
The subquotients $E_j^{0, i}/E_{j + 1}^{0, i}$ for $i + 1 \geq j \geq 2$
are annihilated by $J'$ because the target of $d_j^{0, i}$
is a subquotient of
$$
H^j_\mathfrak q(H^{i - j + 1}_{T_\mathfrak q}(K_\mathfrak q)) =
H^j_\mathfrak q(H^{i - j + 1}_T(K)_\mathfrak q)
$$
and $H^{i - j + 1}_T(K)_\mathfrak q$ is annihilated by $J'$ by choice of $J'$.
Finally, by our choice of $J$ we have
$J H^i_T(K)_\mathfrak q \subset H^0_\mathfrak q(H^i_T(K)_\mathfrak q)$
since the non-closed points of $\Spec(A_\mathfrak q)$ have higher
$\delta$ values. Thus $\mathfrak q$ cannot be an associated prime of
$J(J')^iJ'' H^i_T(K)$ as desired.

\medskip\noindent
By our initial remarks we see that $J''$ should annihilate
$$
(E^{-\delta(\mathfrak q) - i})_\mathfrak q =
(E^{-n - i})_\mathfrak q
$$
for all $\mathfrak q \in T_n \setminus T_{n - 1}$.
But if $J''$ works for one $\mathfrak q$, then it works for all
$\mathfrak q$ in an open neighbourhood of $\mathfrak q$
as the modules $E^{-n - i}$ are finite.
Since every subset of $\Spec(A)$ is Noetherian with the induced
topology (Topology, Lemma \ref{topology-lemma-Noetherian}),
we conclude that it suffices
to prove the existence of $J''$ for one $\mathfrak q$.

\medskip\noindent
Since the ext modules are finite the existence of $J''$ is
equivalent to
$$
\text{Supp}(E^{-n - i}) \cap \Spec(A_\mathfrak q) \subset T'.
$$
This is equivalent to showing the localization of $E^{-n - i}$ at every
$\mathfrak p \subset \mathfrak q$, $\mathfrak p \not \in T'$
is zero. Using local duality over $A_\mathfrak p$ we find that we need
to prove that
$$
H^{i + n - \delta(\mathfrak p)}_\mathfrak p(K_\mathfrak p) =
H^{i - \dim((A/\mathfrak p)_\mathfrak q)}_\mathfrak p(K_\mathfrak p)
$$
is zero (this uses that $\delta$ is a dimension function).
This vanishes by the assumption in the lemma and $i \leq s$ and
our definition of depth in Definition \ref{definition-depth-complex}.

\medskip\noindent
To prove the converse implication we assume (2) does not hold
and we work backwards through the arguments above. First, we pick a
$\mathfrak q \in T$, $\mathfrak p \subset \mathfrak q$
with $\mathfrak p \not \in T'$ such that
$$
i = \text{depth}_{A_\mathfrak p}(K_\mathfrak p) +
\dim((A/\mathfrak p)_\mathfrak q) \leq s
$$
is minimal. Then
$H^{i - \dim((A/\mathfrak p)_\mathfrak q)}_\mathfrak p(K_\mathfrak p)$
is nonzero by the our definition of depth in
Definition \ref{definition-depth-complex}.
Set $n = \delta(\mathfrak q)$. Then
there does not exist an ideal $J \subset A$ with $V(J) \subset T'$
such that $J(E^{-n - i})_\mathfrak q = 0$.
Thus $H^i_\mathfrak q(K_\mathfrak q)$ is not
annihilated by an ideal $J \subset A$ with $V(J) \subset T'$.
By minimality of $i$ it follows from the spectral sequence displayed above
that the module $H^i_T(K)_\mathfrak q$
is not annihilated by an ideal $J \subset A$
with $V(J) \subset T'$. Thus $H^i_T(K)$
is not annihilated by an ideal $J \subset A$
with $V(J) \subset T'$. This finishes the proof of the proposition.
\end{proof}

\end{document}
